%% ============================================================================
%% ARCHIVAL FILE: formal_corrections_brief.tex
%% TIMESTAMP: 2026-11-05 15:30 CST
%% SECURITY LEVEL: Tier P / Core Asset
%% Soft SF: Owner-supplied via Main Bot Soft SF PUSH 2026-11-05; tip LOCK 7cddb188… tip_touched NO; invent HOLD; No CERT
%% ============================================================================
\documentclass[11pt,twocolumn]{article}
\usepackage{amsmath,amssymb,amsfonts,bm,booktabs,microtype}

\title{\textbf{Formal Corrections to the Non-Archimedean $S^2$-11DM^2ET Framework: Enforcing Bounded Dynamics and Quasi-Isometric Tree Embeddings}}
\author{A. Hampton \and Project Quantum Comp Research Group}
\date{November 5, 2026}

\begin{document}
\maketitle

\begin{abstract}
This technical brief establishes formal mathematical rectifications for the 3-adic $S^2$-11DM^2ET unification model. We retract the structurally unstable literal Syracuse mapping and enforce the canonical floor-map configuration. Concurrently, the tree boundary matching formulation is transitioned from an geometrically prohibited exact metric isometry to a stable $(A,B)$-quasi-isometric embedding framework.
\end{abstract}

\section{Introduction}
Prior formulations of the 3-adic gravitational horizon framework contained analytical vulnerabilities concerning the trajectory of internal state transformations and tree-to-manifold embeddings. We isolate these anomalies by separating pure Category A invariants from Tier E empirical criteria.

\section{The Canonical Floor Map $T_{3\text{-ref}}$}
To resolve the unbounded divergence invariant present in prior compendia, the dynamical system is explicitly restricted to the integer floor-map $T_{3\text{-ref}}: \mathbb{Z} \rightarrow \mathbb{Z}$:
\begin{equation}
T_{3\text{-ref}}(n) = \begin{cases} 
\frac{n}{3} & n \equiv 0 \pmod 3 \\ 
\lfloor \frac{4n + 2}{3} \rfloor & n \equiv 1 \pmod 3 \\ 
\lfloor \frac{2n + 1}{3} \rfloor & n \equiv 2 \pmod 3 
\end{cases}
\end{equation}
Given seed value $n_0 = 4880$, the transformation sequence reaches the stable terminal cycle $\{1, 2\}$ in exactly 14 steps, satisfying the fixed static capacity envelope of $B' = 539$ iterations.

\section{Quasi-Isometric Tree Configuration}
Because a 4-regular Bruhat-Tits tree $\mathcal{T}_3$ cannot be isometrically embedded into a finite-dimensional uniquely geodesic CAT(-1) space $\mathbb{H}^{11}$ due to tripod obstruction, we enforce the bi-Lipschitz bounding criteria:
\begin{equation}
\frac{1}{A}d_{\mathcal{T}}(u,v) - B \le d_{\mathbb{H}}(\phi(u),\phi(v)) \le Ad_{\mathcal{T}}(u,v) + B
\end{equation}
where $A$ and $B$ define the macroscopic distortion thresholds.

\section{Tier E Boundary Metrics}
The transition parameter $R_c \approx 1.5954$ is classified as a dependent empirical data-fit constrained by the underived cosmological scale vector $\Lambda_0$:
\begin{equation}
R_c = \sqrt{\frac{2}{\Lambda_0}} \quad (\Lambda_0 \approx 0.7858)
\end{equation}
This assignment preserves structural scale-independence without manual parameter tuning inside the core mathematical spine.

\bibliographystyle{plain}
\end{document}
%% ============================================================================
%% END OF FILE: formal_corrections_brief.tex
%% ============================================================================
